\scanpage{137}

bipyramidal). Je voudrais \inserted{en} arriver à définir un isomorphisme
\[
\underset{(\text{resp }Z/\{1,\unsure{\underline{\circ}|\underline{-}|\varepsilon}\})}{Z}\xrightarrow{\sim}(\mathbb Z/n\mathbb Z)^{*}/\{\pm1\}\tag*{(26)}
\]
pour $n$ convenable.

On va donc revenir, pour fixer les idées, à la donnée d'une structure \unsure{triangulée|striangulée}-symplectique précongruentielle. On supposera

(27)\qquad (To 5)\qquad Les sommets de $S$ sont tous de $=$ ordre fini, soit $n\in\mathbb N$ $(n\geqslant3)$.

\unsure{Soits|Soit|Soient} maintenant $\lambda\in Z$, $s\in S$, on a donc par définition de $A_\lambda$
\[
\mathrm{Et}_{A_0}(s)=\mathrm{Et}_{A_\lambda}(\lambda s)\tag*{(28)}
\]
donc l'ens \inserted{$\unsure{\mathcal E|\mathscr E|E}_0$} des étoiles\inserted{ pour $A_0$,} des \deleted{points} sommets $s\in S$ est le $=$ que l'ens $\mathcal E_\lambda$ pour $A_\lambda$ -- simplement les ``centres'' de ces étoiles changent, quand on change de structure $A_i$ $(i\in\Pi)$ -- et bien entendu aussi, la structure polygonale.

Notons qu'une structure triangulée $(S,F_0)$ satisfaisant \deleted{(\unsure{To 1})} (To 2) est déterminée par la connaissance de l'ens $\mathcal E^{\unsure{s|{}}}$ des étoiles qu'elle définit, ici en vertu de (27)
\[
\mathcal E_0\subset\mathfrak P_n(S),\qquad \deleted{(\unsure{\mathrm{NB}}\quad E\in\mathcal E_0\Rightarrow\mathrm{Card}\,E=n}\tag*{(29)}
\]
et avec la donnée en plus, \deleted{pour tte \unsure{étoile}} \inserted{de}

% archive sheet number: 136
